%---------------------------------------------------------------
%  SECTION II — BACKGROUND AND RELATED WORK
%---------------------------------------------------------------
\section{Background and Related Work}
\label{sec:background}

\subsection{Classical Secure File Transfer}

Traditional SFTP \cite{ylonen2006ssh} and FTPS operate over a
TLS~1.3 \cite{krawczyk2020tls13} or SSH handshake using RSA or
ECDH/ECDSA. These handshakes establish a shared secret in
$O(\log^3 n)$ time classically but are solvable in
$O((\log n)^3)$ time on a quantum computer via Shor's algorithm
\cite{shor1994algorithms}, rendering them fundamentally insecure in
the post-quantum era.

\subsection{Post-Quantum Cryptography}

Post-Quantum Cryptography (PQC) encompasses algorithms believed to
be secure against adversaries equipped with quantum computers
\cite{dam2023survey}. NIST's 2024 standards \cite{nist2024pqc}
formalise three families:

\textbf{CRYSTALS-Kyber (ML-KEM, FIPS~203)} is a Key Encapsulation
Mechanism (KEM) based on the Module-LWE problem. Kyber512 targets
NIST Security Level~1 (equivalent to AES-128 security). Key and
ciphertext sizes are 800~B and 768~B respectively, with key generation
completing in $\approx 2$~ms on commodity hardware
\cite{avanzi2021kyber}.

\textbf{CRYSTALS-Dilithium (ML-DSA, FIPS~204)} is a digital signature
scheme based on Module-LWE and Module-SIS \cite{ducas2018dilithium}.
Dilithium2 targets NIST Security Level~2 with a 1,312-byte public key
and 2,420-byte signature. It provides Existential Unforgeability under
Chosen Message Attack (EUF-CMA).

\textbf{AES-256-GCM} retains 128-bit security against Grover's
quadratic speedup. It provides authenticated encryption (AEAD),
combining CTR-mode encryption with a Galois MAC authentication tag (16
bytes) and a 96-bit nonce (12 bytes), totalling 28 bytes overhead per
encrypted message.

\textbf{BLAKE2b-256} \cite{aumasson2013blake2} offers $\approx 3\times$
the throughput of SHA-256 at equivalent security, making it ideal for
real-time file integrity hashing on both client and server.

\subsection{Related Work}

Sola-Thomas and Imtiaz \cite{sola2025quantumftp} proposed a
quantum-resistant FTP with a blockchain audit trail, validating
Kyber-based key exchange but omitting integrity verification, RBAC,
and a user-facing interface. Bos and Stehl\'{e} \cite{bos2023pqftp}
demonstrated post-quantum FTP using lattice-based key exchange in an
experimental TLS extension but did not address resumable transfers or
metadata privacy. Baseri \cite{baseri2024navigating} surveyed quantum
risks in networked environments but proposed no concrete implementation.
Chen et al.\ \cite{chen2022hybridtls} designed a hybrid PQ-TLS
handshake combining X25519 and Kyber, providing defense-in-depth
against unexpected algorithmic weaknesses.

\textbf{Research gap:} No existing work provides a complete,
deployable post-quantum file transfer system combining PQC handshake,
file integrity, resumable transfers, RBAC, metadata privacy, and a
web GUI in a single containerised application. Q-SFTP fills this gap.
